
\medskip
\noindent\SGARefTarget{sga3:iv:corollary:4:5:4}\textbf{Corollary 4.5.4.}
\textit{Suppose that the given topology on \(\mathscr C\) is coarser than
the canonical topology. For every \(S\in\operatorname{Ob}(\mathscr C)\),
there is an equivalence of categories}
\[
  \widetilde{\mathscr C}_{/\widetilde S}
  \xrightarrow{\ \sim\ }
  \widetilde{\mathscr C/S}.
\]
\textit{The following diagrams commute up to isomorphism; every unlabelled
arrow in them is an equivalence:}
\[
\begin{array}{c@{\;}c@{\;}c@{\;}c@{\;}c@{\;}c@{\;}c}
\mathscr C_{/S}
 &\xrightarrow{(j_{\mathscr C})_{/S}}&
 \widetilde{\mathscr C}_{/\widetilde S}
 &\xrightarrow{(i_{\mathscr C})_{/\widetilde S}}&
 \widehat{\mathscr C}_{/\widehat S}
 &\xrightarrow{(a_{\mathscr C})_{/\widehat S}}&
 \widetilde{\mathscr C}_{/\widetilde S}
 \\
 \big\Vert && \big\downarrow\wr && \big\downarrow\wr && \big\downarrow\wr
 \\
\mathscr C_{/S}
 &\xrightarrow{j_{\mathscr C/S}}&
 \widetilde{\mathscr C/S}
 &\xrightarrow{i_{\mathscr C/S}}&
 \widehat{\mathscr C/S}
 &\xrightarrow{a_{\mathscr C/S}}&
 \widetilde{\mathscr C/S}.
\end{array}
\]
If \(T\to S\) is an object of \(\mathscr C/S\), one also has the
commutative diagram
\[
\begin{array}{ccccc}
 \bigl(\widetilde{\mathscr C}_{/\widetilde S}\bigr)_{/\widetilde T}
 &\longrightarrow&
 \bigl(\widetilde{\mathscr C/S}\bigr)_{/\widetilde T}
 &&\\[-0.15em]
 \Big\downarrow\rlap{\(\scriptstyle\wr\)}
 &&& \searrow &
 \widetilde{(\mathscr C/S)_{/T}}\, .
 \\
 \widetilde{\mathscr C}_{/\widetilde T}
 &\longrightarrow&
 \widetilde{\mathscr C/T}
 &\mathrel{\overset{\sim}{\nearrow}}&
\end{array}
\]

The commutativity of the first two squares follows from the definition of
the equivalence
\(
  \widetilde{\mathscr C}_{/\widetilde S}
  \to\widetilde{\mathscr C/S}
\).
To prove the commutativity of the last square, it is enough to observe that
the following square commutes:
\[
\begin{array}{ccc}
 \widehat{\mathscr C}_{/\widehat S}
 &\xrightarrow{L'}&
 \widehat{\mathscr C}_{/\widehat S}
 \\
 \big\downarrow && \big\downarrow
 \\
 \widehat{\mathscr C/S}
 &\xrightarrow{L''}&
 \widehat{\mathscr C/S}.
\end{array}
\]
Here \(L'\) is the restriction of the functor \(L_{\mathscr C}\) to
\(\widehat{\mathscr C}_{/\widehat S}\), and \(L''\) is the functor
\(L_{\mathscr C/S}\). This is seen immediately by returning to \SGARefLink{sga3:iv:definition:4:3:10:0}{the
definition} of the functors \(L\), given after \SGARefLink{sga3:iv:corollary:4:3:9}{4.3.9}. The second diagram is
simply the restriction to the categories of sheaves of the corresponding
diagram for the categories of presheaves (Expos\'e I, \SGARefLink{sga3:i:section:1}{no.~1}), which is
commutative.

\medskip
\noindent\SGARefTarget{sga3:iv:scholium:4:5:5}\textbf{Scholium 4.5.5.}
The preceding assertions show that, when the given topology is coarser than
the canonical topology, one may identify \(\mathscr C\) with a full
subcategory of \(\widetilde{\mathscr C}\), itself a full subcategory of
\(\widehat{\mathscr C}\). Under these identifications one may make the
customary abuses of language concerning \(\mathscr C\to
\widehat{\mathscr C}\), justified by the preceding commutative diagrams.
In particular, the first diagram of \SGARefLink{sga3:iv:corollary:4:5:4}{4.5.4} shows that the functor \(a\) may
be used without any special precaution.

In the next subsection we shall see that the identification of
\(\mathscr C\) with a full subcategory of \(\widetilde{\mathscr C}\),
unlike the corresponding identification with
\(\widehat{\mathscr C}\), commutes with the formation of certain direct
limits; we shall then explain how this fact is to be used.

\emph{From now on, unless expressly stated otherwise, we assume that the
given topology is coarser than the canonical topology, and we
systematically make the identifications described above.}
